\section{Introduction}
\label{sec:introduction}

A fixed language model can behave very differently when its surrounding
prompts, memory policy, tool interface, or orchestration changes.  This has
motivated systems that optimize prompts and programs
\citep{khattab2023dspy,agrawal2025gepa}, search agent workflows
\citep{hu2025adas}, rewrite full harness code from accumulated traces
\citep{lee2026metaharness,lin2026ahe}, or ask a model to improve its own harness
\citep{zhang2026selfharness}.  These results make automated harness engineering
plausible.  They also make its evaluation unusually adaptive: the same search
process inspects failures, proposes edits, screens candidates, and repeatedly
queries a promotion set.

The object being optimized is therefore not the only source of overfitting.
The \emph{promotion process} is itself a reusable adaptive interface.  A
candidate may score better because of finite-sample noise, a grader shortcut, a
correlated but inactive edit, an unrecorded resource increase, or a regression
outside the optimized slice.  A held-out set helps, but repeatedly consulting
it during search erodes its held-out meaning \citep{dwork2015reusable}.  Even a
replicated score difference leaves a second ambiguity: did the proposed
mechanism produce the change, or did another part of the stochastic execution
do so?

We ask a narrower question than whether a model can rewrite a harness:

\begin{quote}
Given a fixed solver model, mutation space, and search budget, what evidence is
sufficient to promote an adaptively selected harness edit as a reproducible
improvement?
\end{quote}

Our answer is \systemname, a protocol that separates a mutable evolution plane
from an independent verification plane and a one-time sealed authority
(Figure~\ref{fig:mechanism-boundary}).  Each candidate is an immutable child of
one parent and declares a typed, bounded mutation plus a falsifiable mechanism
hypothesis.  Before utility can authorize promotion, matched interventions ask
whether the new code path was activated, whether the agent followed it, and
whether the intended behavior changed.  The gate then requires benefit and
separate non-inferiority constraints for protected behavior and cost.  Missing
or inconsistent evidence cannot be interpreted as a pass.

This design does not assume that mechanism evidence is automatically useful.
It creates a testable treatment contrast.  \fullarm{} exposes typed
activation--adherence--behavior evidence to the optimizer; \scorearm{} receives
only matched aggregate scores, intervals, resource totals, and a frozen
performance-only decision over utility, protected behavior, and cost.  That
projection excludes activation, adherence, and intended-behavior checks; it is
not a redaction of the full mechanism gate.  The solver, proposer, mutation
grammar, task coordinates, and budget are otherwise held fixed.  The real-task
\fullarm{}--\scorearm{} contrast estimates the effect of the complete deployable
policy: it changes both optimizer feedback and the promotion rule.  A
controlled-fault \(2\times2\) factorial separately estimates feedback-guidance,
promotion-verification, and interaction effects.  The central empirical
question is whether the joint \fullarm{} policy reduces false promotion and
improves sealed utility relative to \scorearm{}.

\paragraph{Contributions.}
This draft makes four contributions, with empirical claims deliberately kept
separate from design claims:
\begin{itemize}
  \item A formal promotion problem for adaptive harness search, with typed
  atomic mutations and an explicit trusted surface that candidates cannot edit.
  \item A randomized mechanism gate using child, parent, revert, and matched
  placebo arms, campaign-level error control, and explicit utility, regression,
  safety, and resource constraints.
  \item A prospective factorial design for HarnessFaultBench whose oracle causal
  surface would permit measurement of localization, verified repair, and false
  promotion rather than score alone.
  \item A prospective same-model, budget-matched study across harness surfaces
  and real tasks.  Its protocol and main conclusion are
  \evidenceblocked{pending preregistration, independent search replications, and
  atomic sealed release}.
\end{itemize}

We do not claim the first self-improving harness, activation check, statistical
gate, held-out acceptance test, or use of execution evidence.  The remaining
prospective distinction is child--parent--revert--placebo attribution along an
activation--adherence--behavior chain, campaign-level false-promotion
measurement, and factorial separation of feedback from promotion authority.
